\section{Definitions (RQ2, Part 1)}
\label{sec:definitions}

We build the definitions in strict dependency order, so that none is circular:
syntactic objects first (schema, corpus, deficit), then operational quantities
(reconstruction cost, irrecoverability), then the economic aggregate (debt) and its
decomposition (principal, interest, liability). Each definition ends with its
\emph{observability basis}: what one would measure to instantiate it. The financial
vocabulary is used only \emph{after} the underlying quantity is defined without it.

\begin{definition}[Evidence Schema]
\label{def:schema}
An \emph{evidence schema} $\Sch$ is a declared mapping from each class of
operational change to the required set of \emph{evidence elements} for that class:
typed facts (actor identity, content delta, intent statement, authorization basis,
outcome record, recovery linkage) together with the \emph{link keys} that connect
them across stores. \emph{Observability:} $\Sch$ is an artifact the organization
writes down; its existence is a governance choice, and all subsequent measurement is
relative to it.
\end{definition}

The schema-relativity is deliberate and load-bearing: ``missing'' is undefined
without a statement of what was supposed to exist. This mirrors measurement practice
in data quality~\cite{wang1996} and avoids the circularity of defining debt by
whatever later turns out to be needed. The six-interrogative schema of
\cref{sec:schema-informal} is our reference instantiation; organizations may declare
stricter or weaker ones, and comparisons are meaningful only under a fixed $\Sch$.

\begin{definition}[Evidence Corpus]
\label{def:corpus}
The \emph{evidence corpus} $\Corp(t)$ is the set of evidence elements and link keys
that durably exist and are retrievable at time $t$, across all stores.
\emph{Observability:} enumerable by inventory of the stores, up to discoverability
limits that are themselves part of the phenomenon (\cref{sec:metrics}).
\end{definition}

\begin{definition}[Evidence Deficit]
\label{def:deficit}
For a change $x$ of class $k$ occurring at time $t_x$, the \emph{evidence deficit}
$\Def(x, t) \subseteq \Sch(k)$ is the set of schema-required elements and link keys
for $x$ absent from $\Corp(t)$. The \emph{deficit at creation} is $\Def(x, t_x)$;
elements present at $t_x$ but absent at $t > t_x$ have \emph{decayed}.
\emph{Observability:} a syntactic census---for each change, check each required
element's presence. No judgment about future cost is involved.
\end{definition}

\begin{definition}[Reconstruction and Its Cost]
\label{def:rec}
A \emph{query} $q$ against operational history asks for a set of facts and links
concerning past changes (e.g., ``what changed on service $s$ in window $w$, by whom,
under what authorization?''). \emph{Reconstruction} is the process---human,
algorithmic, or mixed---of attempting to answer $q$ from $\Corp(t)$. Its
\emph{reconstruction-attempt cost} $\ERC(q, \Corp(t))$ is the finite effort expended
(person-time and compute, on a declared accounting) until either (i) an answer of
stated confidence is reached or (ii) the declared policy stops without one. Its
\emph{error} is the divergence of any returned
answer from ground truth. Thus $\ERC$ is defined even when no answer is produced;
policy abstention and certified irrecoverability are represented separately in
\cref{def:debt,def:irrec}. \emph{Observability:} attempt cost is
measurable directly whenever reconstruction is performed (incident retrospectives,
audit preparation, the drills of \cref{sec:field}); error is measurable wherever
ground truth is independently available---natively in simulation
(\cref{sec:simulation}), partially in the field.
A \emph{reconstruction policy} $\rho$ declares the admissible sources, a finite
exhaustion rule, a confidence threshold, and the acceptance/action rule; let
$\mathcal R_q$ be the declared admissible policy class for $q$. The random variable
$C_\rho(q,E)$ is the finite effort used by policy $\rho$ on evidence $E$.
A policy is \emph{sound exhaustive} relative to $(U,\gamma,\rho_{\mathrm{exh}})$
if it completes the declared finite exhaustion and accepts an answer only when the
available evidence establishes it at confidence $\gamma$. Such a policy must
abstain at the certified irrecoverability boundary of \cref{def:irrec}.
Exhaustiveness alone does not imply soundness: an aggressive policy may inspect all
declared sources and still return $W$ after the true fact cannot be established.
Standardized drills estimate policy-specific effort by fixing $\rho$ in advance;
the total-burden optimum used for debt accounting is defined in
\cref{def:debt}, rather than by minimizing effort alone.
\end{definition}

\begin{definition}[Evidence Irrecoverability]
\label{def:irrec}
A fact or link required by query $q$ is \emph{irrecoverable relative to a declared
source universe $U$, confidence threshold $\gamma$, and exhaustion protocol
$\rho_{\mathrm{exh}}$} if no admissible procedure over
$\Corp(t)\cap U$ can establish it at confidence $\gamma$ after
$\rho_{\mathrm{exh}}$ is completed. Write this operational event as
$\Irr(q,t\mid U,\gamma,\rho_{\mathrm{exh}})$. A query is irrecoverable if any fact
essential to it is. This is a scoped certificate, not a metaphysical claim that no
forgotten backup, private message, or offline export exists anywhere.
\emph{Observability:} in simulation, exactly decidable because $U$ and ground truth
are complete; in the field, certified within an enumerated source list, as forensic
practice certifies unavailability~\cite{nist80086,casey2011}.
\end{definition}

\begin{definition}[Evidence Debt]
\label{def:debt}
Fix a schema $\Sch$, a corpus $\Corp(t)$, and a \emph{query workload} $\Qw$: a set
of potential future queries $q$ with occurrence weights $w_q > 0$ (frequencies,
probabilities, or scenario weights, declared explicitly). The default reporting
policy partitions every reconstruction outcome into exactly one of: $S$
(supported-correct), $W$ (accepted-but-wrong), or $N$ (no answer certified by
$\rho$, i.e., policy abstention).
The policy first decides an acceptance indicator $A_\rho(q,E)\in\{0,1\}$ using
only available evidence: $A_\rho=0$ implies $N$ regardless of ground truth; if
$A_\rho=1$, ground truth distinguishes $S$ from $W$. Define
$\Err_q=\{W\}$ and $\mathsf{Abst}_q=\{N\}$, so the accepted-error and policy-
abstention channels are mutually exclusive. Crucially, $\mathsf{Abst}_q$ is not
the irrecoverability event of \cref{def:irrec}: a conservative $\rho$ can abstain
while another admissible procedure could still recover the fact. Write the latter
state separately as $\mathsf I_q(E\mid U,\gamma,\rho_{\mathrm{exh}})\in\{0,1\}$.

For each query template $q$, expectations and probabilities are taken under a
declared evaluation measure $\mu_q$ over query instances and any analyst, policy,
path-survival, or seed randomness. For a deterministic single instance they reduce
to indicators; in the simulation they are empirical frequencies over the
prespecified generated instances and seeds. For a fixed policy, define
\begin{align*}
B_{q,\rho}(E)=&\;\mathbb E C_\rho(q,E)
+\lambda_q\Pr_\rho(\Err_q\mid E)\\[-2pt]
&+\pi_q\Pr_\rho(\mathsf{Abst}_q\mid E),
\end{align*}
and define the optimal total burden
\begin{align*}
B_q^{\star}(E)=\inf_{\rho\in\mathcal R_q}\bigl\{
&\mathbb E C_\rho(q,E)+\lambda_q\Pr_\rho(\Err_q\mid E)\\[-2pt]
&+\pi_q\Pr_\rho(\mathsf{Abst}_q\mid E)\bigr\}.
\end{align*}
The theoretical \emph{optimal evidence debt} of the estate at $t$ is
\begin{equation*}
\ED^{\star}(t)=\sum_{q\in\Qw}w_q
\left[B_q^{\star}(\Corp(t))-B_q^{\star}(\Corp^{*}(t))\right].
\end{equation*}
The policy-specific estimand used by a preregistered drill is instead
\begin{equation*}
\ED_{\rho}(t)=\sum_{q\in\Qw}w_q
\left[B_{q,\rho}(\Corp(t))-B_{q,\rho}(\Corp^{*}(t))\right].
\end{equation*}
Estimating $\ED^{\star}$ requires comparing multiple admissible policies or
declaring $\mathcal R_q$ a singleton; one fixed drill estimates $\ED_\rho$, not an
unobserved infimum. We use unadorned ``ED'' for the construct family and mark every
reported numerical estimate by its policy context.
When the consequence of certified physical loss must differ from a still-
recoverable abstention, replace the additive consequence terms by a declared joint
loss $L_q(Y_\rho,\mathsf I_q)$, where $Y_\rho\in\{S,W,N\}$. This permits, for
example, $L_q(N,0)\ne L_q(N,1)$ and an aggressive wrong answer with
$L_q(W,1)$, without equating a policy decision to a state of the evidence. The
additive instance above is the estimand exercised here: it prices accepted error
and policy-abstention exposure under $\rho$, while IRR is reported separately.

$\Corp^{*}(t)$ is $\Corp(t)$ augmented with every element and link key required by
$\Sch$ for every change up to $t$, captured at change time and subject to no decay.
Thus $\Corp(t)\subseteq\Corp^{*}(t)$. If every policy available on less evidence
can be emulated while ignoring extra evidence, then
$B_q^{\star}(E_2)\leq B_q^{\star}(E_1)$ whenever $E_1\subseteq E_2$; hence
$\ED^{\star}(t)\geq0$. A fixed-policy contrast $\ED_\rho$ inherits this guarantee
only if $\rho$ is information-monotone; the simulation's complete ideal corpus
resolves every link exactly at zero excess burden. This choice is deliberate and consequential: it attributes
retention- and migration-driven loss of once-captured evidence to debt (as
decay-created deficit, \cref{def:deficit}), because from the query's standpoint an
expired record and a never-captured record are the same absence with different
histories. $\lambda_q$ prices an incorrect returned answer and $\pi_q$ prices policy
abstention. Neither coefficient by itself prices certified physical
irrecoverability. \emph{Observability:} estimable, not
directly observable---the estimation problem is exactly what
\cref{sec:metrics,sec:simulation,sec:field} address; the definition fixes what the
estimators are estimators \emph{of}. In field estimation, the same preregistered
policy $\rho$ is applied to actual and ideal/reference corpora to estimate
$\ED_\rho$. In the simulation, the admissible class is the singleton accompanying
reconstructor, so $\ED_\rho=\ED^{\star}$ within that declared class.
\end{definition}

Four features of \cref{def:debt} do argumentative work. First, debt is
\emph{workload-relative}: an estate whose history will never be queried carries
deficits but little debt; declaring $\Qw$ forces the (often implicit) organizational
judgment about which contingencies matter into the open. Second, the excess-cost
form makes the baseline explicit: even a complete corpus makes queries cost
something; debt is only the part deferral added, and optimal information-monotone
total burden makes that excess nonnegative. Third, false reconstruction is priced
explicitly rather than hidden inside effort. Fourth, policy abstention is priced as
a finite consequence, while certified irrecoverability remains a separately
reported state and the mandatory physical boundary of \cref{prop:writeoff}. This
separation prevents a conservative policy from manufacturing ``irrecoverability''
merely by declining to answer.

\begin{definition}[Evidence-Debt Principal]
\label{def:principal}
The \emph{principal} of a deficit $\Def(x, t_x)$ is its contribution to $\ED(t_x)$
at creation time---the excess expected burden were the workload queried immediately,
while memory is fresh, sources are live, and identifiers are current.
\emph{Observability:} estimable at change time from the deficit census and
short-horizon reconstruction calibration.
\end{definition}

\begin{definition}[Evidence-Debt Interest]
\label{def:interest}
The \emph{interest} accrued on change $x$'s evidence position over $[t_x, t]$ is
the growth of its contribution to $\ED$ beyond principal:
$\mathcal{I}(x, t) = \mathrm{ED}\text{-contribution}(x, t) - \text{principal}(x)$.
It has two decay-driven components, both are interest, and they must not be
conflated with new borrowing: (i)~\emph{path decay}, the loss of reconstruction
routes for elements already deficient at $t_x$; and (ii)~\emph{corpus decay}, the
expiry or destruction of elements captured at $t_x$, which enlarges
$\Def(x, t)$ over time per \cref{def:deficit}. The drivers of both are the same
processes: retention expiry, record deletion, personnel
departure~\cite{rigby2016}, resource ephemerality~\cite{burns2016}, identity and
schema drift, vendor migration, and loss of correlation identifiers.
\emph{Observability:} the decay processes are individually observable (retention
policies are published~\cite{cloudtrail,ghdocs,nist80092}; turnover is known to HR;
migrations are scheduled); \cref{sec:interest} shows how their superposition yields
a computable burden-versus-age curve, and why it is not a compound rate.
Interest is gross decay-driven accrual in the absence of remediation. Backfill and
restored linkage are repayments that reduce the debt stock; they are not negative
interest.
\end{definition}

\begin{definition}[Evidence Liability]
\label{def:liability}
The \emph{liability} of the estate with respect to a specific contingency $q$ at $t$
is the conditional realized burden if $q$ occurs at $t$: the actual reconstruction
cost plus the consequence cost of wrongly answered parts or parts left uncertified
by $\rho$ (regulatory
exposure, prolonged outage, lost claims)~\cite{sedona2018,sox2002}.
\emph{Observability:} realized liabilities are observable events (audit findings,
incident duration attributable to archaeology); $\ED$ is their workload-weighted
expectation, which is the standard relation between an expected cost and its
realizations.
\end{definition}

\begin{figure}[t]
\centering
\resizebox{0.70\columnwidth}{!}{% Accounting bridge from evidence deficit to financial exposure (single column)
\begin{tikzpicture}[node distance=2.5mm]
  \node[brok, minimum width=35mm] (deficit)
    {evidence deficit\\{\smlbl missing node, key, or link}};
  \node[stage, below=of deficit, minimum width=35mm] (recon)
    {reconstruction demand\\{\smlbl contingency $q$ activates the gap}};
  \node[human, below=of recon, minimum width=35mm] (hours)
    {engineering effort\\{\smlbl person-hours $+$ compute $+$ delay}};
  \node[evid, below=of hours, minimum width=35mm] (cost)
    {operational / audit cost\\{\smlbl declared rates and delay costs}};
  \node[brok, below=of cost, minimum width=35mm] (exposure)
    {financial exposure\\{\smlbl workload weight $w_q$; loss penalty $\pi_q$}};

  \draw[flow] (deficit) -- (recon);
  \draw[flow] (recon) -- node[right,lbl]{measured ERC} (hours);
  \draw[flow] (hours) -- node[right,lbl]{accounting map} (cost);
  \draw[flow] (cost) -- node[right,lbl]{frequency $\times$ consequence} (exposure);
  \draw[bflow] (recon.west) to[out=190,in=170]
    node[left,lbl,align=right]{irrecoverable\\write-off} (exposure.west);
\end{tikzpicture}
}
\caption{The causal accounting bridge from a syntactic deficit to economic
exposure. Only the first node is obtained by census. Reconstruction drills
measure ERC; declared labor, compute, and delay rates map effort to cost; workload
weights and accepted-error/abstention penalties map policy outcomes to exposure;
physical irrecoverability remains a separate loss state. The figure
defines an auditable chain, not a claim that this paper has estimated dollars.}
\label{fig:economic-chain}
\end{figure}

\Cref{fig:economic-chain} makes explicit the economic bridge that the excess-cost
definition otherwise compresses into one term. A deficit has no intrinsic monetary
value. It becomes economically relevant only when a contingency demands
reconstruction; the resulting person-time, compute, and delay can then be mapped
through published accounting rates, and the workload weights and consequence
penalties of \cref{def:debt} turn those conditional costs into exposure. This
separation is intentional: the paper can define and measure reconstruction burden
without pretending to have measured wages, outage revenue, or regulatory loss.

The resulting vocabulary is internally consistent and independent of the financial
metaphor: \emph{deficit} is syntax (what is missing),
\emph{debt} is expectation (the excess reconstruction effort, accepted error, and
policy abstention the missing evidence is expected to cause),
\emph{principal} is the at-creation component, \emph{interest} is the decay-driven
growth, \emph{liability} is the per-contingency realization, and
\emph{irrecoverability} is the terminal physical loss state, priced separately only
through a declared joint loss when policy abstention is not an adequate proxy.
\Cref{sec:model} gives the machinery that connects them.
